% Options for packages loaded elsewhere
% Options for packages loaded elsewhere
\PassOptionsToPackage{unicode}{hyperref}
\PassOptionsToPackage{hyphens}{url}
\PassOptionsToPackage{dvipsnames,svgnames,x11names}{xcolor}
%
\documentclass[
  ignorenonframetext,
  aspectratio=169,
  spanish,
]{beamer}
\newif\ifbibliography
\usepackage{pgfpages}
\setbeamertemplate{caption}[numbered]
\setbeamertemplate{caption label separator}{: }
\setbeamercolor{caption name}{fg=normal text.fg}
\beamertemplatenavigationsymbolsempty
% remove section numbering
\setbeamertemplate{part page}{
  \centering
  \begin{beamercolorbox}[sep=16pt,center]{part title}
    \usebeamerfont{part title}\insertpart\par
  \end{beamercolorbox}
}
\setbeamertemplate{section page}{
  \centering
  \begin{beamercolorbox}[sep=12pt,center]{section title}
    \usebeamerfont{section title}\insertsection\par
  \end{beamercolorbox}
}
\setbeamertemplate{subsection page}{
  \centering
  \begin{beamercolorbox}[sep=8pt,center]{subsection title}
    \usebeamerfont{subsection title}\insertsubsection\par
  \end{beamercolorbox}
}
% Prevent slide breaks in the middle of a paragraph
\widowpenalties 1 10000
\raggedbottom
\AtBeginPart{
  \frame{\partpage}
}
\AtBeginSection{
  \ifbibliography
  \else
    \frame{\sectionpage}
  \fi
}
\AtBeginSubsection{
  \frame{\subsectionpage}
}
\usepackage{iftex}
\ifPDFTeX
  \usepackage[T1]{fontenc}
  \usepackage[utf8]{inputenc}
  \usepackage{textcomp} % provide euro and other symbols
\else % if luatex or xetex
  \usepackage{unicode-math} % this also loads fontspec
  \defaultfontfeatures{Scale=MatchLowercase}
  \defaultfontfeatures[\rmfamily]{Ligatures=TeX,Scale=1}
\fi
\usepackage{lmodern}

\usetheme[]{metropolis}
\ifPDFTeX\else
  % xetex/luatex font selection
\fi
% Use upquote if available, for straight quotes in verbatim environments
\IfFileExists{upquote.sty}{\usepackage{upquote}}{}
\IfFileExists{microtype.sty}{% use microtype if available
  \usepackage[]{microtype}
  \UseMicrotypeSet[protrusion]{basicmath} % disable protrusion for tt fonts
}{}
\makeatletter
\@ifundefined{KOMAClassName}{% if non-KOMA class
  \IfFileExists{parskip.sty}{%
    \usepackage{parskip}
  }{% else
    \setlength{\parindent}{0pt}
    \setlength{\parskip}{6pt plus 2pt minus 1pt}}
}{% if KOMA class
  \KOMAoptions{parskip=half}}
\makeatother


\usepackage{longtable,booktabs,array}
\usepackage{calc} % for calculating minipage widths
\usepackage{caption}
% Make caption package work with longtable
\makeatletter
\def\fnum@table{\tablename~\thetable}
\makeatother
\usepackage{graphicx}
\makeatletter
\newsavebox\pandoc@box
\newcommand*\pandocbounded[1]{% scales image to fit in text height/width
  \sbox\pandoc@box{#1}%
  \Gscale@div\@tempa{\textheight}{\dimexpr\ht\pandoc@box+\dp\pandoc@box\relax}%
  \Gscale@div\@tempb{\linewidth}{\wd\pandoc@box}%
  \ifdim\@tempb\p@<\@tempa\p@\let\@tempa\@tempb\fi% select the smaller of both
  \ifdim\@tempa\p@<\p@\scalebox{\@tempa}{\usebox\pandoc@box}%
  \else\usebox{\pandoc@box}%
  \fi%
}
% Set default figure placement to htbp
\def\fps@figure{htbp}
\makeatother



\ifLuaTeX
\usepackage[bidi=basic]{babel}
\else
\usepackage[bidi=default]{babel}
\fi
% get rid of language-specific shorthands (see #6817):
\let\LanguageShortHands\languageshorthands
\def\languageshorthands#1{}


\setlength{\emergencystretch}{3em} % prevent overfull lines

\providecommand{\tightlist}{%
  \setlength{\itemsep}{0pt}\setlength{\parskip}{0pt}}



 


\makeatletter
\@ifpackageloaded{tcolorbox}{}{\usepackage[skins,breakable]{tcolorbox}}
\@ifpackageloaded{fontawesome5}{}{\usepackage{fontawesome5}}
\definecolor{quarto-callout-color}{HTML}{909090}
\definecolor{quarto-callout-note-color}{HTML}{0758E5}
\definecolor{quarto-callout-important-color}{HTML}{CC1914}
\definecolor{quarto-callout-warning-color}{HTML}{EB9113}
\definecolor{quarto-callout-tip-color}{HTML}{00A047}
\definecolor{quarto-callout-caution-color}{HTML}{FC5300}
\definecolor{quarto-callout-color-frame}{HTML}{acacac}
\definecolor{quarto-callout-note-color-frame}{HTML}{4582ec}
\definecolor{quarto-callout-important-color-frame}{HTML}{d9534f}
\definecolor{quarto-callout-warning-color-frame}{HTML}{f0ad4e}
\definecolor{quarto-callout-tip-color-frame}{HTML}{02b875}
\definecolor{quarto-callout-caution-color-frame}{HTML}{fd7e14}
\makeatother
\makeatletter
\@ifpackageloaded{caption}{}{\usepackage{caption}}
\AtBeginDocument{%
\ifdefined\contentsname
  \renewcommand*\contentsname{Tabla de contenidos}
\else
  \newcommand\contentsname{Tabla de contenidos}
\fi
\ifdefined\listfigurename
  \renewcommand*\listfigurename{Listado de Figuras}
\else
  \newcommand\listfigurename{Listado de Figuras}
\fi
\ifdefined\listtablename
  \renewcommand*\listtablename{Listado de Tablas}
\else
  \newcommand\listtablename{Listado de Tablas}
\fi
\ifdefined\figurename
  \renewcommand*\figurename{Figura}
\else
  \newcommand\figurename{Figura}
\fi
\ifdefined\tablename
  \renewcommand*\tablename{Tabla}
\else
  \newcommand\tablename{Tabla}
\fi
}
\@ifpackageloaded{float}{}{\usepackage{float}}
\floatstyle{ruled}
\@ifundefined{c@chapter}{\newfloat{codelisting}{h}{lop}}{\newfloat{codelisting}{h}{lop}[chapter]}
\floatname{codelisting}{Listado}
\newcommand*\listoflistings{\listof{codelisting}{Listado de Listados}}
\makeatother
\makeatletter
\makeatother
\makeatletter
\@ifpackageloaded{caption}{}{\usepackage{caption}}
\@ifpackageloaded{subcaption}{}{\usepackage{subcaption}}
\makeatother

\usepackage{bookmark}
\IfFileExists{xurl.sty}{\usepackage{xurl}}{} % add URL line breaks if available
\urlstyle{same}
\hypersetup{
  pdftitle={Aprendizaje y evaluación en la era de la IA},
  pdfauthor={Ph.D.~Pablo Eduardo Caicedo Rodríguez},
  pdflang={es},
  colorlinks=true,
  linkcolor={Maroon},
  filecolor={Maroon},
  citecolor={Blue},
  urlcolor={Blue},
  pdfcreator={LaTeX via pandoc}}


\title{Aprendizaje y evaluación en la era de la IA}
\subtitle{Herramientas, ejemplos y prompts para rediseño evaluativo}
\author{Ph.D.~Pablo Eduardo Caicedo Rodríguez}
\date{2026-07-02}

\begin{document}
\frame{\titlepage}


\section{Tesis central}\label{tesis-central}

\begin{frame}{Tesis central}
\begin{tcolorbox}[enhanced jigsaw, opacityback=0, breakable, left=2mm, bottomrule=.15mm, arc=.35mm, rightrule=.15mm, leftrule=.75mm, toprule=.15mm, colback=white, colframe=quarto-callout-color-frame]

La IA no elimina la evaluación: elimina la validez de las evaluaciones
que solo verifican productos fácilmente automatizables.

\end{tcolorbox}

La tarea institucional es rediseñar el aprendizaje y la evaluación para
certificar \textbf{autonomía}, \textbf{juicio disciplinar},
\textbf{trazabilidad} y \textbf{transferencia}.
\end{frame}

\section{Ruta de la presentación}\label{ruta-de-la-presentaciuxf3n}

\begin{frame}{Ruta de la presentación}
\begin{enumerate}
\tightlist
\item
  Aprendizaje en la era de la IA.
\item
  Evaluación del aprendizaje en la era de la IA.
\item
  Herramientas de evaluación útiles.
\item
  Ejemplos de uso de esas herramientas.
\item
  Prompts para generar instrumentos evaluativos.
\end{enumerate}
\end{frame}

\section{1. Aprendizaje en la era de la
IA}\label{aprendizaje-en-la-era-de-la-ia}

\begin{frame}{Cambio 1: aprender ya no es solo acceder a información}
\phantomsection\label{cambio-1-aprender-ya-no-es-solo-acceder-a-informaciuxf3n}
\begin{tcolorbox}[enhanced jigsaw, rightrule=.15mm, titlerule=0mm, leftrule=.75mm, colframe=quarto-callout-warning-color-frame, opacityback=0, breakable, colbacktitle=quarto-callout-warning-color!10!white, opacitybacktitle=0.6, bottomtitle=1mm, bottomrule=.15mm, toptitle=1mm, arc=.35mm, colback=white, coltitle=black, toprule=.15mm, left=2mm, title=\textcolor{quarto-callout-warning-color}{\faExclamationTriangle}\hspace{0.5em}{Antes\ldots{}}]

Memorizar, recuperar, explicar y entregar.

\end{tcolorbox}

\begin{tcolorbox}[enhanced jigsaw, rightrule=.15mm, titlerule=0mm, leftrule=.75mm, colframe=quarto-callout-caution-color-frame, opacityback=0, breakable, colbacktitle=quarto-callout-caution-color!10!white, opacitybacktitle=0.6, bottomtitle=1mm, bottomrule=.15mm, toptitle=1mm, arc=.35mm, colback=white, coltitle=black, toprule=.15mm, left=2mm, title=\textcolor{quarto-callout-caution-color}{\faFire}\hspace{0.5em}{Ahora\ldots{}}]

Interrogar, verificar, contextualizar, decidir y defender.

\end{tcolorbox}

\begin{tcolorbox}[enhanced jigsaw, rightrule=.15mm, titlerule=0mm, leftrule=.75mm, colframe=quarto-callout-important-color-frame, opacityback=0, breakable, colbacktitle=quarto-callout-important-color!10!white, opacitybacktitle=0.6, bottomtitle=1mm, bottomrule=.15mm, toptitle=1mm, arc=.35mm, colback=white, coltitle=black, toprule=.15mm, left=2mm, title=\textcolor{quarto-callout-important-color}{\faExclamation}\hspace{0.5em}{Riesgo}]

Confundir fluidez generativa con comprensión real.

\end{tcolorbox}

La IA puede apoyar la tutoría, la explicación y la retroalimentación.
\end{frame}

\begin{frame}{Tres niveles de impacto sobre el aprendizaje}
\phantomsection\label{tres-niveles-de-impacto-sobre-el-aprendizaje}
\begin{longtable}[]{@{}
  >{\raggedright\arraybackslash}p{(\linewidth - 4\tabcolsep) * \real{0.3333}}
  >{\raggedright\arraybackslash}p{(\linewidth - 4\tabcolsep) * \real{0.3333}}
  >{\raggedright\arraybackslash}p{(\linewidth - 4\tabcolsep) * \real{0.3333}}@{}}
\toprule\noalign{}
\begin{minipage}[b]{\linewidth}\raggedright
Nivel
\end{minipage} & \begin{minipage}[b]{\linewidth}\raggedright
Qué cambia
\end{minipage} & \begin{minipage}[b]{\linewidth}\raggedright
Pregunta crítica
\end{minipage} \\
\midrule\noalign{}
\endhead
Instrumental & Búsqueda, resumen, edición, traducción & ¿El estudiante
conserva control sobre el problema? \\
Cognitivo & Andamiaje, pistas, feedback, simulaciones & ¿La IA promueve
explicación o reemplaza pensamiento? \\
Epistémico & Qué cuenta como saber, crear o demostrar & ¿La evaluación
mide juicio humano transferible? \\
\bottomrule\noalign{}
\end{longtable}
\end{frame}

\begin{frame}{Rendimiento aparente}
\phantomsection\label{rendimiento-aparente}
\begin{tcolorbox}[enhanced jigsaw, rightrule=.15mm, titlerule=0mm, leftrule=.75mm, colframe=quarto-callout-note-color-frame, opacityback=0, breakable, colbacktitle=quarto-callout-note-color!10!white, opacitybacktitle=0.6, bottomtitle=1mm, bottomrule=.15mm, toptitle=1mm, arc=.35mm, colback=white, coltitle=black, toprule=.15mm, left=2mm, title=\textcolor{quarto-callout-note-color}{\faInfo}\hspace{0.5em}{Nota}]

Una mejora durante el uso de IA no prueba aprendizaje autónomo. Puede
ser solo ejecución asistida.

\end{tcolorbox}

Bastani et al.~muestran que el acceso a IA puede mejorar el desempeño
durante la práctica, pero deteriorar el rendimiento posterior cuando la
herramienta se retira si el sistema entrega soluciones sin andamiaje
pedagógico.
\end{frame}

\begin{frame}{Y el estudiante que \ldots?}
\phantomsection\label{y-el-estudiante-que}
La \textbf{capacidad del estudiante para asumir un papel activo,
intencional y responsable en su propio aprendizaje} no se preserva
prohibiendo IA de forma genérica. Se preserva diseñando tareas donde el
estudiante debe:

\begin{itemize}
\tightlist
\item
  formular hipótesis antes de consultar IA;
\item
  contrastar salidas con fuentes;
\item
  explicar decisiones;
\item
  identificar sesgos, errores o alucinaciones;
\item
  resolver una variante sin asistencia.
\end{itemize}

La asistencia de IA debe evaluarse por su efecto en autorregulación, no
solo por eficiencia.
\end{frame}

\begin{frame}{Principio de diseño del aprendizaje}
\phantomsection\label{principio-de-diseuxf1o-del-aprendizaje}
Una actividad bien diseñada alterna:

\begin{enumerate}
\tightlist
\item
  exploración asistida;
\item
  producción propia;
\item
  crítica de la IA;
\item
  transferencia sin IA;
\item
  defensa argumentada.
\end{enumerate}
\end{frame}

\section{2. Evaluación del aprendizaje en la era de la
IA}\label{evaluaciuxf3n-del-aprendizaje-en-la-era-de-la-ia}

\begin{frame}{La pregunta ya no es ``¿usó IA?''}
\phantomsection\label{la-pregunta-ya-no-es-usuxf3-ia}
La pregunta evaluativa fuerte es:

\begin{quote}
¿Qué evidencia demuestra que el estudiante comprende, decide y
transfiere conocimiento aun cuando la IA participa en el proceso?
\end{quote}

La detección automática de textos generados por IA es insuficiente como
estrategia principal de integridad.
\end{frame}

\begin{frame}{Evaluación vulnerable vs.~evaluación resiliente}
\phantomsection\label{evaluaciuxf3n-vulnerable-vs.-evaluaciuxf3n-resiliente}
\begin{longtable}[]{@{}
  >{\raggedright\arraybackslash}p{(\linewidth - 4\tabcolsep) * \real{0.3333}}
  >{\raggedright\arraybackslash}p{(\linewidth - 4\tabcolsep) * \real{0.3333}}
  >{\raggedright\arraybackslash}p{(\linewidth - 4\tabcolsep) * \real{0.3333}}@{}}
\toprule\noalign{}
\begin{minipage}[b]{\linewidth}\raggedright
Dimensión
\end{minipage} & \begin{minipage}[b]{\linewidth}\raggedright
Vulnerable a IA
\end{minipage} & \begin{minipage}[b]{\linewidth}\raggedright
Resiliente a IA
\end{minipage} \\
\midrule\noalign{}
\endhead
Evidencia & Texto final & Proceso + defensa + transferencia \\
Tarea & Resumir, definir, explicar & Resolver, justificar, auditar,
crear \\
Autoría & Inferida por entrega & Declarada mediante trazabilidad \\
Verificación & Detector o sospecha & Desempeño situado y oralidad \\
Riesgo & Delegación cognitiva & Juicio disciplinar verificable \\
\bottomrule\noalign{}
\end{longtable}
\end{frame}

\begin{frame}{Problema matemático inicial}
\phantomsection\label{problema-matemuxe1tico-inicial}
Un estudiante debe resolver:

\begin{quote}
Una tienda ofrece un descuento del 20\% sobre un producto que cuesta
\$80.000. Luego aplica IVA del 19\% sobre el precio con descuento. ¿Cuál
es el precio final?
\end{quote}
\end{frame}

\begin{frame}{Momento 1: respuesta correcta}
\phantomsection\label{momento-1-respuesta-correcta}
Una IA responde:

\begin{quote}
El descuento es \$16.000. El precio con descuento es \$64.000. El IVA es
\$12.160. El precio final es \$76.160.
\end{quote}

\begin{block}{Pregunta}
\phantomsection\label{pregunta}
¿Esta respuesta demuestra que el estudiante aprendió?
\end{block}
\end{frame}

\begin{frame}{Discusión individual}
\phantomsection\label{discusiuxf3n-individual}
Responda en 3 minutos:

\begin{enumerate}
\tightlist
\item
  ¿Qué parte del procedimiento es correcta?
\item
  ¿Qué evidencia falta para saber si el estudiante comprendió?
\item
  ¿Qué podría haber hecho la IA sin intervención real del estudiante?
\end{enumerate}
\end{frame}

\begin{frame}{Momento 2: rendimiento aparente}
\phantomsection\label{momento-2-rendimiento-aparente}
La respuesta es correcta, pero el estudiante no puede explicar:

\begin{itemize}
\tightlist
\item
  Por qué el IVA se calcula después del descuento.
\item
  Qué representa el 20\%.
\item
  Qué representa el 19\%.
\item
  Cómo verificar si el resultado es razonable.
\end{itemize}

\begin{block}{Pregunta clave}
\phantomsection\label{pregunta-clave}
\begin{quote}
¿Tenemos aprendizaje o solo rendimiento aparente?
\end{quote}
\end{block}
\end{frame}

\begin{frame}{Mini diagnóstico de vulnerabilidad}
\phantomsection\label{mini-diagnuxf3stico-de-vulnerabilidad}
Clasifique las tareas:

\begin{longtable}[]{@{}ll@{}}
\toprule\noalign{}
Tarea & Vulnerabilidad frente a IA \\
\midrule\noalign{}
\endhead
Calcular el precio final & Alta / Media / Baja \\
Explicar cada paso del cálculo & Alta / Media / Baja \\
Resolver una variante sin IA & Alta / Media / Baja \\
Detectar un error en una solución dada & Alta / Media / Baja \\
Defender oralmente el procedimiento & Alta / Media / Baja \\
\bottomrule\noalign{}
\end{longtable}
\end{frame}

\begin{frame}{Trabajo en parejas}
\phantomsection\label{trabajo-en-parejas}
Analicen esta solución alternativa:

\begin{quote}
Precio inicial: \$80.000 IVA: \$15.200 Precio con IVA: \$95.200
Descuento del 20\%: \$19.040 Precio final: \$76.160
\end{quote}

\begin{block}{Preguntas}
\phantomsection\label{preguntas}
\begin{enumerate}
\tightlist
\item
  ¿El resultado final coincide?
\item
  ¿El procedimiento es equivalente?
\item
  ¿Qué diferencia hay entre descontar antes del IVA y descontar después?
\item
  ¿Qué debería justificar el estudiante?
\end{enumerate}
\end{block}
\end{frame}

\begin{frame}{Variación del problema}
\phantomsection\label{variaciuxf3n-del-problema}
Ahora resuelva sin IA:

\begin{quote}
Un producto cuesta \$120.000. Primero se aplica un descuento del 15\%.
Luego se cobra IVA del 19\%. ¿Cuál es el precio final?
\end{quote}

\begin{block}{Evidencia requerida}
\phantomsection\label{evidencia-requerida}
No basta con entregar el número. Debe mostrar:

\begin{itemize}
\tightlist
\item
  Procedimiento.
\item
  Justificación.
\item
  Verificación del resultado.
\item
  Interpretación económica.
\end{itemize}
\end{block}
\end{frame}

\begin{frame}{Bitácora mínima de trazabilidad}
\phantomsection\label{bituxe1cora-muxednima-de-trazabilidad}
Si se permite IA, el estudiante debe declarar:

\begin{longtable}[]{@{}ll@{}}
\toprule\noalign{}
Elemento & Evidencia \\
\midrule\noalign{}
\endhead
Prompt usado & ¿Qué le pidió a la IA? \\
Respuesta recibida & ¿Qué obtuvo? \\
Verificación & ¿Cómo comprobó el cálculo? \\
Corrección & ¿Detectó errores o ambigüedades? \\
Aprendizaje & ¿Qué entendió después del proceso? \\
\bottomrule\noalign{}
\end{longtable}
\end{frame}

\begin{frame}{Defensa breve}
\phantomsection\label{defensa-breve}
Cada estudiante responde oralmente:

\begin{enumerate}
\tightlist
\item
  ¿Por qué el descuento se calcula antes del IVA?
\item
  ¿Cómo comprobaría que el resultado es razonable?
\item
  ¿Qué cambiaría si el IVA se aplicara antes del descuento?
\item
  ¿Qué parte del ejercicio no debería delegarse a la IA?
\end{enumerate}
\end{frame}

\begin{frame}{Cierre conceptual}
\phantomsection\label{cierre-conceptual}
Una evaluación es vulnerable cuando solo exige:

\begin{quote}
Resultado correcto.
\end{quote}

Una evaluación es más resiliente cuando exige:

\begin{quote}
Procedimiento, explicación, verificación, transferencia y defensa.
\end{quote}
\end{frame}

\begin{frame}{Transición hacia AIAS}
\phantomsection\label{transiciuxf3n-hacia-aias}
Este ejercicio permite introducir una escala de uso de IA:

\begin{longtable}[]{@{}ll@{}}
\toprule\noalign{}
Nivel & Uso de IA \\
\midrule\noalign{}
\endhead
0 & IA prohibida \\
1 & IA solo para comprender instrucciones \\
2 & IA para apoyo inicial \\
3 & IA para resolver con trazabilidad \\
4 & IA para comparar soluciones \\
5 & IA como objeto de auditoría crítica \\
\bottomrule\noalign{}
\end{longtable}
\end{frame}

\begin{frame}{Pregunta final}
\phantomsection\label{pregunta-final}
\begin{quote}
En este problema matemático, ¿en qué nivel permitiría usted el uso de
IA?
\end{quote}
\end{frame}

\begin{frame}{Para tener en cuenta al rediseñar la evaluación}
\phantomsection\label{para-tener-en-cuenta-al-rediseuxf1ar-la-evaluaciuxf3n}
\begin{itemize}
\tightlist
\item
  Validez de la evaluación.
\item
  Autonomía del estudiante.
\item
  Trazabilidad del proceso.
\item
  Juicio humano.
\item
  Transferencia a un nuevo problema.
\end{itemize}
\end{frame}

\begin{frame}{La evaluación formativa gana centralidad}
\phantomsection\label{la-evaluaciuxf3n-formativa-gana-centralidad}
La literatura clásica ya indicaba que la retroalimentación eficaz reduce
la brecha entre desempeño actual y desempeño esperado.

En la era de la IA, la retroalimentación debe cumplir cuatro
condiciones:

\begin{itemize}
\tightlist
\item
  ser técnicamente correcta;
\item
  ser accionable;
\item
  exigir reflexión metacognitiva;
\item
  preparar para transferencia.
\end{itemize}
\end{frame}

\begin{frame}{AIAS: declarar niveles de uso de IA}
\phantomsection\label{aias-declarar-niveles-de-uso-de-ia}
La Artificial Intelligence Assessment Scale permite comunicar con
precisión qué nivel de IA es aceptable en una evaluación.

\begin{longtable}[]{@{}rll@{}}
\toprule\noalign{}
Nivel & Uso de IA & Evidencia requerida \\
\midrule\noalign{}
\endhead
0 & Sin IA & Desempeño individual controlado \\
1 & Planeación mínima & Declaración simple de uso \\
2 & Apoyo limitado & Registro de consultas y decisiones \\
3 & Colaboración estructurada & Bitácora, crítica y edición humana \\
4 & Uso amplio & Auditoría de salidas y defensa \\
5 & Exploración abierta & Reflexión ética, técnica y creativa \\
\bottomrule\noalign{}
\end{longtable}
\end{frame}

\begin{frame}{Regla de oro para rediseñar}
\phantomsection\label{regla-de-oro-para-rediseuxf1ar}
\begin{tcolorbox}[enhanced jigsaw, opacityback=0, breakable, left=2mm, bottomrule=.15mm, arc=.35mm, rightrule=.15mm, leftrule=.75mm, toprule=.15mm, colback=white, colframe=quarto-callout-color-frame]

Si una IA puede resolver la tarea con calificación alta sin interacción
significativa del estudiante, la tarea no está evaluando bien.

\end{tcolorbox}

La respuesta no es únicamente sancionar; es rediseñar la evidencia de
aprendizaje.
\end{frame}

\begin{frame}{Criterios mínimos de validez}
\phantomsection\label{criterios-muxednimos-de-validez}
\begin{longtable}[]{@{}ll@{}}
\toprule\noalign{}
Criterio & Evidencia observable \\
\midrule\noalign{}
\endhead
Autonomía & Resuelve una variante sin IA \\
Trazabilidad & Explica prompts, fuentes, descartes y decisiones \\
Juicio & Detecta errores de IA y justifica correcciones \\
Transferencia & Aplica el conocimiento en contexto nuevo \\
Ética & Declara uso, límites, datos y responsabilidad \\
\bottomrule\noalign{}
\end{longtable}
\end{frame}

\section{3. Herramientas de evaluación
útiles}\label{herramientas-de-evaluaciuxf3n-uxfatiles}

\begin{frame}{Herramienta 1: portafolio con trazabilidad}
\phantomsection\label{herramienta-1-portafolio-con-trazabilidad}
El portafolio deja de ser un repositorio de productos. Debe contener:

\begin{itemize}
\tightlist
\item
  versiones sucesivas;
\item
  prompts usados;
\item
  fuentes verificadas;
\item
  decisiones de aceptación/rechazo;
\item
  reflexión metacognitiva;
\item
  evidencia final defendible.
\end{itemize}

Uso adecuado: cursos de proyecto, investigación, programación, análisis
de datos, diseño o escritura académica.
\end{frame}

\begin{frame}{Herramienta 2: defensa oral breve}
\phantomsection\label{herramienta-2-defensa-oral-breve}
La defensa oral permite verificar comprensión y autoría cognitiva.

\begin{longtable}[]{@{}ll@{}}
\toprule\noalign{}
Tipo de pregunta & Qué verifica \\
\midrule\noalign{}
\endhead
``Explique por qué descartó esta salida'' & Juicio crítico \\
``Resuelva una variante del caso'' & Transferencia \\
``Defienda su rúbrica'' & Criterio disciplinar \\
``Identifique un error de la IA'' & Alfabetización crítica \\
\bottomrule\noalign{}
\end{longtable}
\end{frame}

\begin{frame}{Herramienta 3: simulaciones y desempeño en vivo}
\phantomsection\label{herramienta-3-simulaciones-y-desempeuxf1o-en-vivo}
Las simulaciones desplazan la evaluación del texto a la acción.

Ejemplos:

\begin{itemize}
\tightlist
\item
  OSCE/OSKI clínico. \textbf{Objective Structured Clinical
  Examination}/\textbf{Objective Structured Knowledge/Skill Interview}.
\item
  Laboratorio con procedimiento observado.
\item
  Comité técnico simulado.
\item
  Sustentación de modelo de datos.
\item
  Diagnóstico de caso con evidencia incompleta.
\end{itemize}
\end{frame}

\begin{frame}{Herramienta 4: rúbricas analíticas asistidas por IA}
\phantomsection\label{herramienta-4-ruxfabricas-analuxedticas-asistidas-por-ia}
La IA puede apoyar el diseño, pero no debe sustituir el juicio docente.

\begin{longtable}[]{@{}lll@{}}
\toprule\noalign{}
Fase & Rol de IA & Rol docente \\
\midrule\noalign{}
\endhead
Borrador & Propone criterios & Ajusta a resultados de aprendizaje \\
Niveles & Sugiere descriptores & Elimina ambigüedad y sesgos \\
Prueba & Simula respuestas & Calibra con ejemplos reales \\
Uso & Ayuda a retroalimentar & Decide calificación final \\
\bottomrule\noalign{}
\end{longtable}
\end{frame}

\begin{frame}{Herramienta 5: banco de preguntas generadas y auditadas}
\phantomsection\label{herramienta-5-banco-de-preguntas-generadas-y-auditadas}
La IA puede generar variantes, casos y distractores. El proceso válido
exige:

\begin{enumerate}
\tightlist
\item
  alineación con resultado de aprendizaje;
\item
  revisión disciplinar;
\item
  control de dificultad;
\item
  identificación de respuesta esperada;
\item
  explicación de distractores;
\item
  prueba piloto.
\end{enumerate}
\end{frame}

\begin{frame}{Herramienta 6: revisión por pares con criterios
explícitos}
\phantomsection\label{herramienta-6-revisiuxf3n-por-pares-con-criterios-expluxedcitos}
La revisión por pares es más sólida cuando la IA se usa para:

\begin{itemize}
\tightlist
\item
  sugerir preguntas de revisión;
\item
  detectar vacíos argumentativos;
\item
  comparar contra rúbrica;
\item
  generar retroalimentación preliminar.
\end{itemize}

El estudiante debe aceptar, rechazar o corregir cada sugerencia.
\end{frame}

\begin{frame}{Herramienta 7: auditoría de salidas de IA}
\phantomsection\label{herramienta-7-auditoruxeda-de-salidas-de-ia}
Evaluar al estudiante como auditor:

\begin{longtable}[]{@{}ll@{}}
\toprule\noalign{}
Producto de IA & Tarea del estudiante \\
\midrule\noalign{}
\endhead
Resumen & Verificar omisiones y distorsiones \\
Código & Probar, depurar, explicar complejidad \\
Diagnóstico & Identificar supuestos y riesgos \\
Rúbrica & Detectar sesgos y criterios vagos \\
Bibliografía & Confirmar DOI, pertinencia y existencia \\
\bottomrule\noalign{}
\end{longtable}
\end{frame}

\begin{frame}{Herramienta 8: matriz de alineación}
\phantomsection\label{herramienta-8-matriz-de-alineaciuxf3n}
La matriz evita tareas tecnológicamente vistosas pero pedagógicamente
débiles.

\begin{longtable}[]{@{}
  >{\raggedright\arraybackslash}p{(\linewidth - 8\tabcolsep) * \real{0.1875}}
  >{\raggedright\arraybackslash}p{(\linewidth - 8\tabcolsep) * \real{0.1875}}
  >{\raggedleft\arraybackslash}p{(\linewidth - 8\tabcolsep) * \real{0.2500}}
  >{\raggedright\arraybackslash}p{(\linewidth - 8\tabcolsep) * \real{0.1875}}
  >{\raggedright\arraybackslash}p{(\linewidth - 8\tabcolsep) * \real{0.1875}}@{}}
\toprule\noalign{}
\begin{minipage}[b]{\linewidth}\raggedright
Resultado de aprendizaje
\end{minipage} & \begin{minipage}[b]{\linewidth}\raggedright
Evidencia
\end{minipage} & \begin{minipage}[b]{\linewidth}\raggedleft
Nivel AIAS
\end{minipage} & \begin{minipage}[b]{\linewidth}\raggedright
Instrumento
\end{minipage} & \begin{minipage}[b]{\linewidth}\raggedright
Verificación
\end{minipage} \\
\midrule\noalign{}
\endhead
Analizar & Dictamen & 3 & Caso situado & Defensa oral \\
Crear & Prototipo & 4 & Portafolio & Demostración \\
Evaluar & Auditoría & 4 & Revisión crítica & Variante sin IA \\
Aplicar & Procedimiento & 0/2 & Simulación & Observación \\
\bottomrule\noalign{}
\end{longtable}
\end{frame}

\section{4. Ejemplos de uso}\label{ejemplos-de-uso}

\begin{frame}{Ejemplo 1: ensayo académico rediseñado}
\phantomsection\label{ejemplo-1-ensayo-acaduxe9mico-rediseuxf1ado}
\begin{longtable}[]{@{}
  >{\raggedright\arraybackslash}p{(\linewidth - 2\tabcolsep) * \real{0.5000}}
  >{\raggedright\arraybackslash}p{(\linewidth - 2\tabcolsep) * \real{0.5000}}@{}}
\toprule\noalign{}
\begin{minipage}[b]{\linewidth}\raggedright
Antes
\end{minipage} & \begin{minipage}[b]{\linewidth}\raggedright
Después
\end{minipage} \\
\midrule\noalign{}
\endhead
Ensayo final sobre un tema & Dossier de fuentes + mapa argumental +
borradores + defensa \\
Calificación por texto & Calificación por proceso, evidencia y juicio \\
Prohibición genérica de IA & AIAS 3: IA permitida con trazabilidad \\
\bottomrule\noalign{}
\end{longtable}

\textbf{Evidencia clave:} el estudiante debe defender por qué una fuente
fue incluida, excluida o corregida.
\end{frame}

\begin{frame}{Ejemplo 2: proyecto de ciencia de datos}
\phantomsection\label{ejemplo-2-proyecto-de-ciencia-de-datos}
Secuencia evaluativa:

\begin{enumerate}
\tightlist
\item
  Formulación del problema sin IA.
\item
  Exploración de enfoques con IA declarada.
\item
  Construcción del pipeline.
\item
  Auditoría de sesgos y errores.
\item
  Reproducción de resultados.
\item
  Defensa técnica con cambios en vivo.
\end{enumerate}

\textbf{Verificación:} modificación de datos o métrica durante la
defensa.
\end{frame}

\begin{frame}{Ejemplo 3: programación}
\phantomsection\label{ejemplo-3-programaciuxf3n}
\begin{longtable}[]{@{}ll@{}}
\toprule\noalign{}
Evidencia & Criterio \\
\midrule\noalign{}
\endhead
Repositorio con commits & Evolución real del trabajo \\
Prompts y respuestas & Uso transparente de IA \\
Pruebas unitarias & Funcionamiento verificable \\
Explicación de errores & Comprensión del código \\
Refactorización en vivo & Autonomía técnica \\
\bottomrule\noalign{}
\end{longtable}

La evaluación no premia solo que el código corra; premia que el
estudiante pueda explicarlo, corregirlo y transferirlo.
\end{frame}

\begin{frame}{Ejemplo 4: evaluación clínica o biomédica}
\phantomsection\label{ejemplo-4-evaluaciuxf3n-cluxednica-o-biomuxe9dica}
Instrumento: simulación tipo OSCE/OSKI.

\begin{itemize}
\tightlist
\item
  IA permitida para preparar un protocolo.
\item
  Sin IA durante el desempeño observado.
\item
  Preguntas orales sobre riesgos, supuestos y decisiones.
\item
  Rúbrica con criterios técnicos, éticos y comunicativos.
\end{itemize}

\textbf{Resultado:} la IA apoya la preparación; la competencia se
verifica en ejecución.
\end{frame}

\begin{frame}{Ejemplo 5: evaluación de fuentes con DOI}
\phantomsection\label{ejemplo-5-evaluaciuxf3n-de-fuentes-con-doi}
Tarea:

\begin{enumerate}
\tightlist
\item
  La IA propone 10 fuentes.
\item
  El estudiante verifica existencia, DOI y pertinencia.
\item
  Clasifica evidencia: revisión, experimento, marco conceptual o
  política.
\item
  Detecta fuentes inexistentes o mal citadas.
\item
  Presenta una matriz de trazabilidad.
\end{enumerate}

\textbf{Competencia:} alfabetización informacional y juicio académico.
\end{frame}

\begin{frame}{Ejemplo 6: debate técnico con dossier}
\phantomsection\label{ejemplo-6-debate-tuxe9cnico-con-dossier}
\begin{longtable}[]{@{}lll@{}}
\toprule\noalign{}
Fase & Uso de IA & Evaluación \\
\midrule\noalign{}
\endhead
Preparación & Sí, declarado & Calidad del dossier \\
Debate & No o limitado & Argumentación y contraargumento \\
Dictamen & Sí, para revisión & Fundamentación final \\
Defensa & No & Dominio conceptual \\
\bottomrule\noalign{}
\end{longtable}
\end{frame}

\section{5. Prompts útiles para generar herramientas de
evaluación}\label{prompts-uxfatiles-para-generar-herramientas-de-evaluaciuxf3n}

\begin{frame}{Prompt 1: alinear evaluación con resultados de
aprendizaje}
\phantomsection\label{prompt-1-alinear-evaluaciuxf3n-con-resultados-de-aprendizaje}
\begin{quote}
Actúa como diseñador curricular universitario. A partir del siguiente
resultado de aprendizaje: {[}pegar resultado{]}, diseña una evaluación
resiliente al uso de IA generativa. Entrega: objetivo evaluativo,
evidencia esperada, nivel AIAS recomendado, criterios de logro, forma de
verificar autonomía y mecanismo de trazabilidad. Evita tareas que puedan
resolverse solo con producción textual automatizada.
\end{quote}
\end{frame}

\begin{frame}{Prompt 2: crear rúbrica analítica}
\phantomsection\label{prompt-2-crear-ruxfabrica-analuxedtica}
\begin{quote}
Diseña una rúbrica analítica para evaluar {[}actividad{]}. Debe tener 4
criterios, 4 niveles de desempeño y descriptores observables. Incluye un
criterio de trazabilidad del uso de IA y otro de transferencia sin IA.
No uses criterios vagos como ``buena calidad''; cada descriptor debe
indicar evidencia observable.
\end{quote}
\end{frame}

\begin{frame}{Prompt 3: convertir una tarea vulnerable en resiliente}
\phantomsection\label{prompt-3-convertir-una-tarea-vulnerable-en-resiliente}
\begin{quote}
Evalúa la vulnerabilidad a IA de esta tarea: {[}pegar tarea{]}.
Identifica qué partes puede resolver una IA sin aprendizaje real.
Rediseña la actividad para incluir proceso, bitácora, defensa oral,
transferencia y verificación de fuentes. Propón un nivel AIAS y explica
la razón pedagógica.
\end{quote}
\end{frame}

\begin{frame}{Prompt 4: diseñar defensa oral}
\phantomsection\label{prompt-4-diseuxf1ar-defensa-oral}
\begin{quote}
Genera 12 preguntas de defensa oral para verificar si el estudiante
comprende realmente este trabajo: {[}pegar resumen del trabajo{]}.
Organiza las preguntas en: comprensión conceptual, decisiones
metodológicas, crítica de IA, transferencia a un caso nuevo y ética.
Incluye una respuesta esperada breve para cada pregunta.
\end{quote}
\end{frame}

\begin{frame}{Prompt 5: generar banco de preguntas con validación}
\phantomsection\label{prompt-5-generar-banco-de-preguntas-con-validaciuxf3n}
\begin{quote}
Crea un banco de 15 preguntas para evaluar {[}tema{]}. Incluye 5
preguntas de aplicación, 5 de análisis y 5 de evaluación. Para cada
pregunta entrega: resultado de aprendizaje asociado, respuesta esperada,
distractores justificados si es de selección múltiple, nivel de
dificultad y advertencia sobre posible respuesta generada por IA.
\end{quote}
\end{frame}

\begin{frame}{Prompt 6: construir una bitácora de IA}
\phantomsection\label{prompt-6-construir-una-bituxe1cora-de-ia}
\begin{quote}
Diseña una plantilla de bitácora para documentar el uso de IA en una
evaluación. Debe incluir: propósito de cada consulta, prompt usado,
respuesta obtenida, decisión del estudiante, fuente externa usada para
verificar, cambio realizado en el producto y reflexión sobre qué se
aprendió.
\end{quote}
\end{frame}

\begin{frame}{Prompt 7: auditar una respuesta de IA}
\phantomsection\label{prompt-7-auditar-una-respuesta-de-ia}
\begin{quote}
Actúa como evaluador crítico. Revisa la siguiente respuesta generada por
IA: {[}pegar respuesta{]}. Identifica errores conceptuales, afirmaciones
sin fuente, sesgos, omisiones y riesgos éticos. Luego propone cómo un
estudiante debería corregirla y qué evidencia debería presentar para
demostrar aprendizaje.
\end{quote}
\end{frame}

\begin{frame}{Prompt 8: diseñar evaluación por simulación}
\phantomsection\label{prompt-8-diseuxf1ar-evaluaciuxf3n-por-simulaciuxf3n}
\begin{quote}
Diseña una simulación evaluativa para verificar la competencia
{[}competencia{]}. Define escenario, información inicial, restricciones,
acciones esperadas del estudiante, criterios observables, errores
críticos, preguntas de defensa y forma de calificación. Indica qué uso
de IA se permite antes, durante y después.
\end{quote}
\end{frame}

\begin{frame}{Prompt 9: crear matriz de alineación}
\phantomsection\label{prompt-9-crear-matriz-de-alineaciuxf3n}
\begin{quote}
A partir de este programa de curso: {[}pegar programa{]}, construye una
matriz con columnas: resultado de aprendizaje, actividad, evidencia,
nivel AIAS, herramienta de evaluación, criterios de rúbrica, mecanismo
de trazabilidad y verificación sin IA. Señala incoherencias entre
resultados, tareas y evaluación.
\end{quote}
\end{frame}

\begin{frame}{Prompt 10: generar retroalimentación formativa}
\phantomsection\label{prompt-10-generar-retroalimentaciuxf3n-formativa}
\begin{quote}
Genera retroalimentación formativa para este trabajo: {[}pegar
trabajo{]}. No asignes calificación. Organiza el feedback en:
fortalezas, brechas frente a la rúbrica, preguntas de mejora, acciones
concretas y una tarea corta de transferencia. Evita reescribir el
trabajo por el estudiante.
\end{quote}
\end{frame}

\section{Síntesis operativa}\label{suxedntesis-operativa}

\begin{frame}{Síntesis operativa}
Para evaluar en la era de la IA:

\begin{enumerate}
\tightlist
\item
  Defina el nivel de IA permitido antes de entregar la tarea.
\item
  Evalúe proceso, no solo producto.
\item
  Exija trazabilidad verificable.
\item
  Incorpore defensa, simulación o transferencia.
\item
  Use IA para apoyar diseño y feedback, no para delegar el juicio final.
\end{enumerate}
\end{frame}

\section{Lista rápida de instrumentos
recomendados}\label{lista-ruxe1pida-de-instrumentos-recomendados}

\begin{frame}{Lista rápida de instrumentos recomendados}
\begin{longtable}[]{@{}lll@{}}
\toprule\noalign{}
Objetivo & Instrumento principal & Complemento \\
\midrule\noalign{}
\endhead
Comprensión profunda & Defensa oral & Preguntas de transferencia \\
Proyecto complejo & Portafolio & Bitácora de IA \\
Juicio crítico & Auditoría de salida IA & Matriz de fuentes \\
Competencia práctica & Simulación & Observación con rúbrica \\
Escritura académica & Dossier + versiones & Verificación DOI \\
Programación & Repositorio + pruebas & Refactorización en vivo \\
\bottomrule\noalign{}
\end{longtable}
\end{frame}

\section{Cierre}\label{cierre}

\begin{frame}{Cierre}
\begin{quote}
La evaluación válida en la era de la IA no pregunta solo quién escribió
el texto. Pregunta quién comprende, quién decide, quién verifica y quién
puede transferir el conocimiento cuando la asistencia desaparece.
\end{quote}
\end{frame}




\end{document}
